\section{Introduction}
\label{sec:intro}

Tax administrations that monitor enrolled debt need a forecast of cash that will be collected, not only a description of the stock.
Enrolled debt is tax debt that has been formally registered for collection.
The stock is large and persistent.
Realized collection is seasonal.
At a daily resolution it is also sparse: weekends, public holidays, and days without a payment slip are structural zeros, not missing measurements that should be dropped before scoring.
A claim that one model is best is meaningful only for a stated grain and a stated metric.

\paragraph{Research question.}
Under one evaluation protocol, which family among the Stage~1 forecasters M0--M9 and the Stage~2 energy and generative forecasters E1--E4 best predicts collection on the monthly stock-and-collection panel and on the daily and weekly payment-slip aggregates, and what does that ranking imply for a monitoring instrument that may abstain?

\paragraph{Protocol in brief.}
We forecast four series drawn from the same extract.
The monthly series is a panel of enrolled-debt stock and collection.
The calendar-daily series (Ranking~A) is the sum of payment-slip values on a continuous civil calendar, with explicit zeros.
The weekly series sums that daily total by week ending Sunday.
A fourth series, business-day daily (Ranking~B), keeps Monday--Friday observations and is an ablation of the daily target.
Brazilian national holidays stay in Ranking~B.
They are not removed.
All families share contiguous train, validation, and test blocks, validation sMAPE for selection, and test sMAPE for the official rank.

\paragraph{Contributions.}
\begin{itemize}
  \item A multi-grain forecasting protocol for enrolled-debt collection, with official ranking by raw test sMAPE and a complementary skill score against the seasonal naive of the same grain.
  \item A Stage~2 comparison that keeps selective TEM (an energy model with an abstention threshold) distinct from a disagreement proxy that only switches between a boosting forecast and the seasonal naive.
  \item Evidence, on the March~2026 extract, that the sMAPE leader is grain-specific, and that retaining calendar zeros versus restricting the series to business days changes both the error level and the ordering.
\end{itemize}

Section~\ref{sec:related} places the protocol among classical forecasting, deep time-series models, and selective prediction.
Section~\ref{sec:data} defines the series.
Section~\ref{sec:method} states metrics, selection, and each model class.
Section~\ref{sec:results} reports the freeze.
Sections~\ref{sec:discussion}--\ref{sec:limitations} interpret the ranking and state what the extract does not support.
